\begin{figure*}
    \centering
    \captionsetup{justification=raggedright,singlelinecheck=false}
    \includegraphics[width=1\linewidth]{1_overview.pdf}
    \caption{Mid Training of recent models. Only selected models with detailed data mix are shown. In the long context part, Hybrid means a combination of different attention mechanisms, e.g., self-attention, SSM, lightning-attention, etc. INT is the abbreviation for interleaving global and local attentions. }
    \label{fig:data_lr_longcontext}
\end{figure*}

\section{Introduction}
\label{sec:intro}


In recent large language model (LLM) training, mid-training has emerged as a crucial continuation stage between general pre-training and task-specific supervised fine-tuning (SFT). Rather than relying solely on web-scale corpora, which provide broad but noisy supervision, state-of-the-art models increasingly adopt an annealing-style phase to refine optimization dynamics and sharpen data quality~\cite{dubey2024llama, hu2024minicpm, glorioso2024zamba}. This stage is motivated by both practical and theoretical considerations: lowering the learning rate mitigates gradient variance and stabilizes convergence near favorable minima, while shifting toward more curated or synthetic corpora~\cite{abdin2024phi3} enhances the marginal utility of each additional token. Without such adjustments, models often stagnate in capability development, showing limited gains in reasoning, coding, and long-context understanding despite substantial additional compute~\cite{bi2024deepseek, allal2025smollm2}. 

The effectiveness of mid-training can be theoretically motivated, as it shifts models from memorization toward abstraction by emphasizing structured reasoning, factuality, and instruction following. We summarize this phenomenon from three perspectives: (i) gradient noise scale, where high-quality data improves optimization by enhancing signal variance and mitigating overfitting plateaus~\cite{mccandlish2018empirical, wang2021eliminating, meng2023per}; (ii) the information bottleneck, which compresses noisy features while preserving predictive structures~\cite{izmailov2022feature}; and (iii) curriculum learning, where data distributions are gradually refined to reinforce complex reasoning~\cite{wang2025dump, jia2022sample}. Together, these perspectives provide a principled rationale for mid-training as a strategy to improve generalization and efficiency in large-scale training.

Despite its growing adoption and effectiveness, there has been no comprehensive survey that conceptualizes mid-training as a coherent paradigm. Existing studies have explored isolated aspects (e.g., curriculum-style data annealing~\cite{du2024survey,wu2025data}, adaptive learning-rate schedules~\cite{jin2023rethinking}, or context-length extension~\cite{huang2023advancing, liu2025comprehensive, pawar2024and}) but lack a comprehensive framework. To bridge this gap, we present an integrated review of mid-training strategies across three interconnected domains: data distribution, optimization scheduling, and context extension. By framing these elements as mutually reinforcing, we highlight mid-training as a distinct and coherent stage in the LLM development pipeline, rather than an ad-hoc collection of heuristics.

A central aspect of mid-training is the refinement of the data mixture after large-scale pre-training. At this point, models have acquired broad linguistic and semantic features from trillions of web-scale tokens, making the marginal value of additional noisy data limited. To improve efficiency, the distribution is shifted toward curated or synthetic corpora emphasizing reasoning, coding, STEM, multilinguality, and instruction following. Empirical evidence shows that such targeted data can foster compositional skills (e.g., chain-of-thought reasoning, math problem solving, and multi-step planning) that are underrepresented in general crawls~\cite{abdin2024phi, allal2025smollm2}. Common practice includes down-sampling low-quality tokens and up-sampling knowledge-dense corpora, often through iterative learning rate annealing passes~\cite{NEURIPS2024_19e4ea30, penedo2024fineweb, Su2024NemotronCCTC}. 

Another key component of mid-training is adapting the learning rate schedule. After aggressive pre-training, models approach promising optima where large step sizes risk instability or divergence. Reducing the learning rate with smoother decay stabilizes convergence and suppresses gradient noise~\cite{you1908does, gilmer2022loss, kalra2024warmup}. Existing strategies include linear or cosine decay, multi-stage schedulers, and adaptive schemes tailored to long runs~\cite{goyal2017accurate, loshchilov2016sgdr, bi2024deepseek, hu2024minicpm, shen2024power, ibrahim2024simple}. These methods improve sample efficiency by enabling finer assimilation of high-quality tokens. Yet scheduler design remains largely empirical: the optimal duration and shape of decay vary across model sizes, architectures, and optimizers~\cite{kaplan2020scaling, hoffmann2022training, defazio2023optimal, bergsma2025straight, li2025predictable}. Overly steep decay may halt learning prematurely, while overly cautious schedules waste computation with limited gains.

Finally, mid-training is also a natural stage for extending context length beyond the 4K–8K limits of early pre-training. This extension supports downstream tasks that demand coherence across long documents, multi-file reasoning, or dialogue histories spanning tens of thousands of tokens. Current methods combine positional encoding remapping, such as Position Interpolation~\cite{chen2023extending}, NTK-aware interpolation~\cite{bloc97_ntk_rope_2023}, YaRN~\cite{peng2023yarn} or ABF~\cite{xiong2023effective} with curricula that gradually introduce long-form inputs. Studies show that extending context during mid-training enables models to capture long-range dependencies more effectively without restarting pre-training. 

Together, these elements form the backbone of mid-training: without refined data, the model lacks specialization; without proper scheduling, it risks instability or underutilization of valuable tokens; and without context extension, it remains constrained in scope for downstream applications. Their interdependence suggests that effective mid-training cannot be understood by analyzing each factor in isolation, but rather through their combined impact on model performance and efficiency. This holistic view motivates the need for systematic discussion of mid-training as a distinct stage in the LLM development pipeline.

\textbf{Our contributions}. This paper contributes to academia and industry in the following ways:
\begin{itemize}[leftmargin=10pt]
    \item We propose a taxonomy of LLM mid-training based on three domains (data distribution, learning rate scheduling and long-context extension), and summarize various approaches in each domain. To the best of our knowledge, it is the first such taxonomy on LLM mid-training.
    \item We summarize valuable insights in each domain, providing readers with a convenient reference for future use.
    \item We discuss the common evaluation benchmarks and collect the reported gains from different models, thereby, providing a structured comparison of how mid-training contribute to LLM training.
    \item We outline promising future research direction in LLM mid-training and propose potential ways forward.
\end{itemize}
